Um einen genaueren Überblick über die Entwicklung von \ac{ERP}-Systemen zu erhalten zu können, ist es hilfreich, die historischen Meilensteine nachzuvollziehen, die zur Entstehung und Weiterentwicklung dieser Systeme beigetragen haben.

Die Ursprünge der heutigen \ac{ERP}-Systeme reichen zurück bis in die 1960er Jahre. Zu dieser Zeit begangen Unternehmen erste simple computergestützte Lösungen für das Bedarfskontrolle zu nutzen, welche dazu dienten Anforderungen an den Bestand zu identifizieren und dessen Nutzung zu überwachen. Zu dieser frühen Phase wurde die Rechenleistung hauptsächlich für algorithmische Berechnungen verwendet. \cite[\pagef 2ff]{Bititci_Strategic_1998}